
\begin{ledgroupsized}[r]{120mm}
\footnotesize
\pstart
\noindent\textbf{Überlieferung:}
\pend
\end{ledgroupsized}
\begin{ledgroupsized}[r]{114mm}
\footnotesize
\pstart
\parindent -6mm
\makebox[6mm][l]{\textit{L}}%
Aufzeichnung:
LH~XXXV~10, 14~Bl.~3. 
Ein Zettel (ca.~10,5~x~10~cm.);		%
rechter und unterer Rand schräg abgeschnitten.
Zwei Seiten; Textfolge durch den Verweis "\textit{verte}" auf Bl.~3~r\textsuperscript{o} gesichert.
In der oberen linken Ecke von Bl.~3~r\textsuperscript{o}, in Rot, der Vermerk von unklarer Hand \textit{:\;III}
(vgl. die Träger der benachbarten Signaturen LH~XXXV~10, 14 Bl.~2 bis 13, u.\,a. 
N.\;\ref{60117},
\ref{60118},
\ref{60115},
\ref{60112},
\ref{60116},
\ref{60113},
\ref{60114},
\ref{41824}
und \ref{60111}
in diesem Band).
\pend
\end{ledgroupsized}
\vspace{5mm}
\begin{ledgroup}
\footnotesize
\pstart
\noindent
\textbf{Datierungsgründe:}
Die Datierung der vorliegenden Aufzeichnung N.\;\ref{60110} beruht auf der Erwähnung des spätestens im Mai 1692 verfassten, um 1693 veröffentlichten Aufsatzes 
\cite{00542}\textit{Regle generale de la composition des mouvemens} (\cite{00157}\textit{JS} vom 7.~September 1693, S.\;417\textendash419; \cite{01043}\textit{LMG} VI, S.\;231--233)
bzw. der zwei "problèmes", die Leibniz auf den Aufsatz folgen ließ:
\cite{00545}\textit{Deux problemes construits par M.~Leibniz} (ebd. 14.~September 1693, S.\;423\,f.; \cite{01043}\textit{LMG} VI, S.\;233\,f.; die \textit{Regle} und die \textit{Problemes} erscheinen vsl. in einem künftigen Band von \textit{LSB} VIII). 
In letzterer Publikation behandelt Leibniz die Zerlegung der "conatus" eines Punktes, der von verschiedenen Seilen in ebensoviele Richtungen gezogen wird – ein Ansatz, den er in der vorliegenden Aufzeichnung N.\;\ref{60110} auf den Fall einer Kugel, die auf von Seilen gestützten Hebeln liegt, anzuwenden sucht. 
(Zur Anwendung der \textit{Regle generale} von 1692/1693 in der Statik, siehe auch das Konzept "Quae sit proportio gravitatis ponderis ex pluribus chordis suspensi ad pondera sustinentia": N.\;\ref{60105}.)
\pend 
\end{ledgroup}
\vspace*{8mm}
\selectlanguage{latin}
\pstart
\noindent
\href{\Scan{LH_35_10_14_003r}}%
{\lbrack3~r\textsuperscript{o}\rbrack}
\pend
\nobreak
\begin{edfigure}%
\hfill
\includegraphics[scale=1]{%
gesamttex/edit_VIII-4/images/LH_35_10_14_003_d1_003r.pdf}
\hspace*{2mm}%
\hfill%
\includegraphics[scale=1]{%
gesamttex/edit_VIII-4/images/LH_35_10_14_003_d2_003r.pdf}
\hfill
\includegraphics[scale=1]{%
gesamttex/edit_VIII-4/images/LH_35_10_14_003_d3_003r.pdf}
\hfill
\edcaption{%
\hspace*{1mm}%
\lbrack\textit{Fig.~1}\rbrack 
\hspace*{20mm}%
\lbrack\textit{Fig.~2}\rbrack
\hspace*{58mm}%
\lbrack\textit{Fig.~3}\rbrack
\hspace*{42mm}%
}
\end{edfigure}
\pstart
\noindent
Angulo \textit{DCE} cujus latera \textit{DC}, \textit{EC}, sint
\edtext{mobilia singulatim}{%
\lemma{mobilia}%
\Bfootnote{%
\textit{(1)}~circa %
\textit{(2)}~singulatim~\textit{L}}}
circa centrum fixum \textit{C} 
\edtext{incumbat\indexS{globus duobus planis incumbens} 
globus \textit{G} in punctis \textit{D} et \textit{E} tangens,}{%
\lemma{incumbat}%
\Bfootnote{%
\textit{(1)}~in punctis \textit{D} et \textit{E} globus
\textit{(2)}~globus \bracks{...} tangens,~\textit{L}}}
quaeritur quantum latera premat,\indexS{pressio globi in plana}
seu quanta debeant esse pondera \textit{A} et \textit{B} quibus laterum \textit{DC}, \textit{EC} conversio impeditur. %
\pend
\pstart
Per
\edtext{Globi\indexS{globus} \textit{G} centrum ducantur recta}{%
\lemma{Globi \textit{G}}%
\Bfootnote{%
\textit{(1)}~punctum %
\textit{(2)}~centrum ducantur
\textit{(a)}~rectae
\textit{(b)}~recta~\textit{L}}}
\textit{GM} normalis ad
\edtext{\textit{CD}, exprimens conatum\indexS{conatus ponderis} ponderis \textit{A} in \textit{G}, et}{%
\lemma{\textit{CD},}%
\Bfootnote{%
\textit{(1)}~et %
\textit{(2)}~exprimens %
\textit{(a)}~vis %
\textit{(b)}~conatum ponderis \textit{A}
\bars{%
\textit{(1)}~in \textit{G}
\textit{(2)}~in \textit{G}}, 
et~\textit{L}}}
recta \textit{GN}, normalis ad \textit{CE}, exprimens conatum\indexS{conatus ponderis} ponderis \textit{B} in \textit{G}.
Denique compleatur rectangulum \textit{GNLM} erit pondus \textit{G} ut
\edtext{diagonalis \textit{GL}, et pondera}{%
\lemma{diagonalis \textit{GL},}%
\Bfootnote{%
\textit{(1)}~et \textit{A}, \textit{B} eru %
\textit{(2)}~et pondera~\textit{L}}}
\textit{A}, \textit{B}, erunt ut
\edtext{latera \textit{GM}, \textit{GN}. Debet}{%
\lemma{latera \textit{GM}, \textit{GN}.}%
\Bfootnote{%
\textit{(1)}~Intelligi %
\textit{(2)}~Debet~\textit{L}}}
autem intelligi brachia quibus pondera aguntur \textit{DH} et \textit{EK} item \textit{HA}, \textit{KB} esse aequalia. %
\pend
\begin{edfigure}%
\includegraphics[scale=1]{%
gesamttex/edit_VIII-4/images/LH_35_10_14_003_d4_003r.pdf}
\edcaption{\lbrack\textit{Fig.~4}\rbrack}
\end{edfigure}


\pstart
Scilicet perinde est, ac si
\edtext{pondus\indexS{pondus trahens globum fune}
\textit{(A)} ut \textit{GM}}{%
\lemma{pondus \textit{(A)}}%
\Bfootnote{%
\textit{(1)}~traheret %
\textit{(2)}~ut \textit{GM}~\textit{L}}}
traheret fune \textit{(A)GM}, et pondus\indexS{pondus trahens globum fune}
\textit{(B)} ut \textit{GN} traheret fune
\edlabel{35_10_14_003v_1a}%
\edtext{}{%
\xxref{35_10_14_003v_1a}{35_10_14_003v_1b}%
\lemma{\textit{(B)GN}.}%
\Bfootnote{%
\href{\Scan{LH_35_10_14_003v}}%
{\lbrack3~v\textsuperscript{o}\rbrack}\ 
\textit{(1)}~Quaecunque %
\textit{(2)}~Quotcunque~\textit{L}}}%
\textit{(B)GN}.
\pend
\pstart
\href{\Scan{LH_35_10_14_003v}}%
{\lbrack3~v\textsuperscript{o}\rbrack}\ 
Quotcunque\indexS{globus pluribus planis incumbens}%
\edlabel{35_10_14_003v_1b} 
sint latera planorum,\indexS{planum inclinatum} et quoscunque angulos faciant, res erit reducta ad totidem funes, et expedietur 
\edtext{ex mea regula\indexS{regula generalis compositionis motuum} generali compositionis
\edtext{mot\bracks{uu}m.}{%
\lemma{}%
\Bfootnote{motum \textit{L ändert Hrsg.}}}%
}{%
\lemma{ex mea regula generali compositionis mot\bracks{uu}m}%
\Cfootnote{\indexN{\textls{Leibniz} (Leibnitius, GGL), Gottfried Wilhelm 1646\textendash1716}\textsc{G.\,W.~Leibniz},
\cite{00542}\textit{Regle generale de la composition des mouvemens}, \cite{00157}\textit{JS} vom 7.~September 1693, S.\;417\textendash419 und \textsc{Ders.}, \cite{00545}\textit{Deux problemes construits par M.~Leibniz}, ebd. 14.~September 1693, S.\;423\,f. (\cite{01043}\textit{LMG} VI, S.\;231--233 und 233\,f.).}} 
Nempe pondera
\edtext{talia esse}{%
\lemma{talia}%
\Bfootnote{%
\textit{(1)}~essent %
\textit{(2)}~esse~\textit{L}}}
debent, ut directiones quas gravi \textit{G} imprimunt, exprimendo rectis pro ponderum magnitudine ex \textit{G} productis,
\edtext{et quarum}{%
\lemma{et}%
\Bfootnote{%
\textit{(1)}~punctorum loca directionum %
\textit{(2)}~quarum~\textit{L}}}
extrema erunt loca directionum. Horum centrum gravitatis\indexS{centrum gravitatis}
\edtext{inventum jungatur}{%
\lemma{inventum}%
\Bfootnote{%
\textit{(1)}~debet %
\textit{(2)}~jungatur~\textit{L}}}
ipsi \textit{G} centro\indexS{centrum gravitatis} Gravis, recta haec jungens debet esse verticalis; eademque multipla per numerum ponderum debet esse proportionalis ipsi ponderi \textit{G}. Hinc
\edtext{patet cum duo\indexS{globus duobus planis incumbens}
sunt plana}{%
\lemma{patet}%
\Bfootnote{%
\textit{(1)}~in duobus ponde %
\textit{(2)}~cum duo sunt plana~\textit{L}}}
rem nos consequi uno tantum modo, sed idem debet fieri et cum pluribus\indexS{globus pluribus planis incumbens}
planis, utique enim determinatum est quantum prematur.
\pend
\pstart
Quid si globus impellatur non in
\edtext{puncto\indexS{globus in puncto impellens} sed}{%
\lemma{puncto}%
\Bfootnote{%
\textit{(1)}~sed in plano tunc quasi infiniti erunt imp %
\textit{(2)}~sed %
\textit{(3)}~sed~\textit{L}}}
in superficie,\indexS{globus in superficie, non in puncto impellens} 
tunc erunt\indexS{conatus impressi infiniti}
infiniti impressi conatus. Et pondus\indexS{pondus in superficie impellens} 
superficie impellens fingi potest in omnia ejus puncta distributum. 
\pend
